% !PRINTED none
% [sic] no printed folio number visible on this section-opening page (Tenth Section); page contains the section title, the \S.\,108 heading, and running body text, but no running head/number is legible in any of the three bands.

\begin{center}
\gesperrt{Tenth Section.}

\textbf{Approximate Calculation of the Roots.}

\rule{3cm}{0.4pt}

\S.\ 108.

\gesperrt{Interpolation. Regula falsi.}
\end{center}

With the delimitation of the intervals in each of which only one root of an algebraic equation lies, the possibility is given of calculating the real roots numerically with arbitrary accuracy, by narrowing the interval more and more, e.g.\ by repeated halving. If one divides an interval of size 1 repeatedly into ten parts, then each new step yields a further decimal place.

Once a root has first been separated from all the others, then in the further narrowing of the interval one has only to take into account the sign of the function $f(x)$ itself. The calculations needed for this can be greatly facilitated by the application of the interpolation formula which we came to know in the first section.

Namely, in \S.\,10, (7) we derived a formula by which one can determine a function of degree $n$, which we shall now denote by $\varphi(\xi)$, when $n+1$ successive values $\varphi(0), \varphi(1) \ldots \varphi(n)$ are given. This formula is, if
\weq{1}{\begin{aligned}
B_\nu^{(\xi)} &= \frac{\xi(\xi-1)\ldots(\xi-\nu+1)}{1.2.3\ldots\nu}\\
\Delta_\xi &= \varphi(\xi+1) - \varphi(\xi),
\end{aligned}}
%CONT>>
